\chapter{2001年大纲卷（文）}
\section{单选题}
\begin{problem}
\(\tan 300^{\circ} + \cot 405^{\circ}\) 的值为（\quad）

\choices
    {\(1 + \sqrt{3}\)}
    {\(1 - \sqrt{3}\)}
    {\(-1 - \sqrt{3}\)}
    {\(-1 + \sqrt{3}\)}
\end{problem}

\begin{answer}
B
\end{answer}

\begin{solution}
\(300^{\circ}=360^{\circ}-60^{\circ}\)，\(405^{\circ}=360^{\circ}+45^{\circ}\)，所以
\(\tan 300^{\circ}=-\tan 60^{\circ}=-\sqrt{3}\)，\(\cot 405^{\circ}=\cot 45^{\circ}=1\). 因而原式为 \(1-\sqrt{3}\)，对应选项 B.
\end{solution}

\begin{problem}
过点 \(A(1, -1)\)，\(B(-1, 1)\) 且圆心在直线 \(x + y - 2 = 0\) 上的圆的方程是（\quad）

\choices
    {\((x - 3)^2 + (y + 1)^2 = 4\)}
    {\((x + 3)^2 + (y - 1)^2 = 4\)}
    {\((x - 1)^2 + (y - 1)^2 = 4\)}
    {\((x + 1)^2 + (y + 1)^2 = 4\)}
\end{problem}

\begin{answer}
C
\end{answer}

\begin{solution}
设圆心为 \(O(u,v)\). 因为圆经过 \(A\)，\(B\)，所以 \(OA=OB\)，即
\((u-1)^2+(v+1)^2=(u+1)^2+(v-1)^2\). 展开并整理得 \(u=v\). 又因为 \(O\) 在直线 \(x+y-2=0\) 上，所以 \(u+v=2\)，从而 \(u=v=1\)，圆心为 \(O(1,1)\). 此时 \(r^2=OA^2=(-1)^2+(-2)^2=4\)，故圆的方程为 \((x-1)^2+(y-1)^2=4\)，对应选项 C.
\end{solution}

\begin{problem}
若一个圆锥的轴截面积是等边三角形，其面积为 \(\sqrt{3}\)，则这个圆锥的全面积是（\quad）

\choices
    {\(3\pi\)}
    {\(3\sqrt{3}\pi\)}
    {\(6\pi\)}
    {\(9\pi\)}
\end{problem}

\begin{answer}
A
\end{answer}

\begin{solution}
设圆锥的底面半径为 \(r\)，母线长为 \(l\). 轴截面是以底面直径为底、以两条母线为腰的等腰三角形，因此其三边均相等时有 \(l=2r\). 轴截面面积为 \(\sqrt{3}\)，故
\(\frac{\sqrt{3}}{4}(2r)^2=\sqrt{3}\)，从而 \(r=1\)，\(l=2\). 圆锥的侧面积为 \(\pi rl=2\pi\)，底面积为 \(\pi r^2=\pi\)，所以全面积为 \(3\pi\)，对应选项 A.
\end{solution}

\begin{problem}
若定义在区间 \((-1,0)\) 内的函数 \(f(x) = \log_{2a}(x + 1)\) 满足 \(f(x) > 0\)，则 \(a\) 的取值范围是（\quad）

\choices
    {\(\left(0, \frac{1}{2}\right)\)}
    {\(\left(0, \frac{1}{2}\right]\)}
    {\(\left(\frac{1}{2},+\infty\right)\)}
    {\((0, +\infty)\)}
\end{problem}

\begin{answer}
A
\end{answer}

\begin{solution}
对数有意义要求底数满足 \(2a>0\) 且 \(2a\ne1\). 当 \(x\in(-1,0)\) 时，真数 \(x+1\in(0,1)\). 对任意 \(t\in(0,1)\)，当且仅当底数满足 \(0<2a<1\) 时有 \(\log_{2a}t>0\)：若 \(2a>1\)，则 \(\log_{2a}t<0\)，而 \(2a=1\) 时对数无意义. 因而 \(0<a<\frac12\)，对应选项 A.
\end{solution}

\begin{problem}
已知复数 \(z = \sqrt{2} + \sqrt{6}\i\)，则 \(\arg \frac{1}{z}\) 是（\quad）

\choices
    {\(\frac{\pi}{3}\)}
    {\(\frac{5\pi}{3}\)}
    {\(\frac{\pi}{6}\)}
    {\(\frac{11\pi}{6}\)}
\end{problem}

\begin{answer}
B
\end{answer}

\begin{solution}
复数 \(z=\sqrt{2}+\sqrt{6}\i\) 的实部、虚部均为正，故其辐角在第一象限. 又
\(\tan\arg z=\frac{\sqrt{6}}{\sqrt{2}}=\sqrt{3}\)，所以 \(\arg z=\frac{\pi}{3}\). 取倒数只改变辐角的符号，故 \(\arg\frac{1}{z}=-\frac{\pi}{3}\pmod{2\pi}\). 按选项采用 \([0,2\pi)\) 内的辐角表示，得到 \(\arg\frac{1}{z}=\frac{5\pi}{3}\)，对应选项 B.
\end{solution}

\begin{problem}
函数 \(y = 2^{-x} + 1\)（\(x > 0\)）的反函数是（\quad）

\choices
    {\( y = \log_2\frac{1}{x - 1} \)（\(x \in (1,2)\)）}
    {\(y = -\log_2\frac{1}{x - 1}\)（\(x \in (1,2)\)）}
    {\( y = \log_2\frac{1}{x - 1} \)（\(x \in (1,2]\)）}
    {\( y = -\log_2\frac{1}{x - 1} \)（\(x \in (1,2]\)）}
\end{problem}

\begin{answer}
A
\end{answer}

\begin{solution}
当原函数的自变量满足 \(x>0\) 时，\(0<2^{-x}<1\)，所以其值域为 \(1<y<2\). 交换自变量与因变量，得到 \(x=2^{-y}+1\)，即 \(x-1=2^{-y}\). 因此
\(y=-\log_2(x-1)=\log_2\frac{1}{x-1}\)，且反函数的定义域为原函数的值域 \(x\in(1,2)\). 这与选项 A 一致.
\end{solution}

\begin{problem}
若椭圆经过原点，且焦点为 \(F_{1}(1,0)\)，\(F_{2}(3,0)\)，则其离心率为（\quad）

\choices
    {\(\frac{3}{4}\)}
    {\(\frac{2}{3}\)}
    {\(\frac{1}{2}\)}
    {\(\frac{1}{4}\)}
\end{problem}

\begin{answer}
C
\end{answer}

\begin{solution}
椭圆中心是两焦点的中点 \((2,0)\)，焦半距为 \(c=1\). 原点在椭圆上，且原点到两个焦点的距离分别为 \(1\)，\(3\)，由椭圆定义有
\(2a=1+3=4\)，所以 \(a=2\). 因而离心率 \(e=\frac{c}{a}=\frac{1}{2}\)，对应选项 C.
\end{solution}

\begin{problem}
若 \(0 < \alpha < \beta < \frac{\pi}{4}\)，\(\sin \alpha + \cos \alpha = a\)，\(\sin \beta + \cos \beta = b\)，则（\quad）

\choices
    {\(a > b\)}
    {\(a < b\)}
    {\(ab < 1\)}
    {\(ab > 2\)}
\end{problem}

\begin{answer}
B
\end{answer}

\begin{solution}
设 \(h(t)=\sin t+\cos t\)，则 \(h'(t)=\cos t-\sin t\). 在 \(0<t<\frac{\pi}{4}\) 时有 \(\cos t>\sin t\)，所以 \(h(t)\) 在该区间上严格递增. 由 \(\alpha<\beta\) 得 \(a=h(\alpha)<h(\beta)=b\). 另外 \(h(0)=1\)，\(h\left(\frac{\pi}{4}\right)=\sqrt{2}\)，故 \(1<a<b<\sqrt{2}\)，从而 \(ab>1\) 且 \(ab<2\)，选项 \(ab<1\) 与 \(ab>2\) 均不成立. 因此正确选项为 B.
\end{solution}

\begin{problem}
在正三棱柱 \(ABC - A_{1}B_{1}C_{1}\) 中，若 \(AB = \sqrt{2} BB_{1}\)，则 \(AB_{1}\) 与 \(C_{1}B\) 所成的角的大小为（\quad）

\choices
    {\(60^{\circ}\)}
    {\(90^{\circ}\)}
    {\(45^{\circ}\)}
    {\(120^{\circ}\)}
\end{problem}

\begin{answer}
B
\end{answer}

\begin{solution}
设 \(AB=s\)，则棱柱高 \(h=BB_1=\frac{s}{\sqrt{2}}\). 取
\(A=(0,0,0)\)，\(B=(s,0,0)\)，\(C=\left(\frac{s}{2},\frac{\sqrt{3}s}{2},0\right)\)，于是
\(B_1=(s,0,h)\)，\(C_1=\left(\frac{s}{2},\frac{\sqrt{3}s}{2},h\right)\). 可取两直线的方向向量
\(\overrightarrow{AB_1}=(s,0,h)\)，\(\overrightarrow{C_1B}=\left(\frac{s}{2},-\frac{\sqrt{3}s}{2},-h\right)\). 它们的内积为
\(\overrightarrow{AB_1}\cdot\overrightarrow{C_1B}=\frac{s^2}{2}-h^2=\frac{s^2}{2}-\frac{s^2}{2}=0\). 因而 \(AB_1\perp C_1B\)，所成角为 \(90^{\circ}\)，对应选项 B.
\end{solution}

\begin{problem}
设 \(f(x)\)，\(g(x)\) 都是单调函数，有如下四个命题中，正确的命题是

\begin{circlelist}
\item 若 \(f(x)\) 单调递增，\(g(x)\) 单调递增，则 \(f(x) - g(x)\) 单调递增；
\item 若 \(f(x)\) 单调递增，\(g(x)\) 单调递减，则 \(f(x) - g(x)\) 单调递增；
\item 若 \(f(x)\) 单调递减，\(g(x)\) 单调递增，则 \(f(x) - g(x)\) 单调递减；
\item 若 \(f(x)\) 单调递减，\(g(x)\) 单调递减，则 \(f(x) - g(x)\) 单调递减.
\end{circlelist}

（\quad）

\choices
    {\(\circled{1}\circled{3}\)}
    {\(\circled{1}\circled{4}\)}
    {\(\circled{2}\circled{3}\)}
    {\(\circled{2}\circled{4}\)}
\end{problem}

\begin{answer}
C
\end{answer}

\begin{solution}
若 \(x_1<x_2\)，且 \(f\) 单调递增、\(g\) 单调递减，则 \(f(x_1)\leqslant f(x_2)\)，\(g(x_1)\geqslant g(x_2)\)，相减得到
\(f(x_1)-g(x_1)\leqslant f(x_2)-g(x_2)\)，所以命题\(\circled{2}\)正确. 同理，若 \(f\) 单调递减、\(g\) 单调递增，则
\(f(x_1)-g(x_1)\geqslant f(x_2)-g(x_2)\)，所以命题\(\circled{3}\)正确.

命题\(\circled{1}\)不一定成立，例如取 \(f(x)=x\)，\(g(x)=2x\)，两者都递增，但 \(f(x)-g(x)=-x\) 递减. 命题\(\circled{4}\)也不一定成立，例如取 \(f(x)=-x\)，\(g(x)=-2x\)，两者都递减，但 \(f(x)-g(x)=x\) 递增. 因此只有命题\(\circled{2}\)，\(\circled{3}\)正确，对应选项 C.
\end{solution}

\begin{problem}
一间平房的屋顶有如图三种不同的盖法，分别是：

\begin{circlelist}
\item 单向倾斜；
\item 双向倾斜；
\item 四向倾斜.
\end{circlelist}

记三种盖法屋顶面积分别为 \(P_{1}\)，\(P_{2}\)，\(P_{3}\). 若屋顶斜面与水平面所成的角都是 \(\alpha\)，则（\quad）

\problemresetlayout
\bitmapfigure[max width=0.74\linewidth]{img/2001/syllabus_liberal/q11_fig1.png}

\choices
    {\(P_{3} > P_{2} > P_{1}\)}
    {\(P_{3} > P_{2} = P_{1}\)}
    {\(P_{3} = P_{2} > P_{1}\)}
    {\(P_{3} = P_{2} = P_{1}\)}
\end{problem}

\begin{answer}
D
\end{answer}

\begin{solution}
设平房屋顶在水平面上的投影面积为 \(S_0\). 对任一种盖法，将各个斜面面积分别投影到水平面上；由于每个斜面与水平面所成的角均为 \(\alpha\)，总投影面积满足
\(S_0=P_i\cos\alpha\)，其中 \(i=1\)，\(2\)，\(3\). 三种盖法覆盖的是同一屋顶，故水平投影面积相同；又 \(\cos\alpha>0\)，所以 \(P_1=P_2=P_3\). 因此对应选项 D.
\end{solution}

\begin{problem}
如图，小圆圈表示网络的结点，结点之间的连线表示它们有网线相联，连线标注的数字表示该段网线单位时间内可以通过的最大信息量，现从结点 \(A\) 向结点 \(B\) 传递信息，信息可以分开沿不同的路线同时传递，则单位时间内传递的最大信息量为（\quad）

\bitmapfigure[max width=0.42\linewidth]{img/2001/syllabus_liberal/q12_fig1.png}

\choices
    {\(26\)}
    {\(24\)}
    {\(20\)}
    {\(19\)}
\end{problem}

\begin{answer}
D
\end{answer}

\begin{solution}
把图中上方两条路线的汇合结点记为 \(U\)，下方两条路线的汇合结点记为 \(V\). 上方信息从 \(B\) 出发时必须经过容量分别为 \(3\)，\(4\) 的两条网线，因此上方路线的总流量不超过 \(3+4=7\). 下方信息在到达 \(A\) 前必须共同经过容量为 \(12\) 的网线，因此下方路线的总流量不超过 \(12\). 两组路线在图中只在端点汇合，所以总流量不超过 \(7+12=19\).

取上方两条路线分别传递 \(3\) 和 \(4\) 个单位，下方两条路线分别传递 \(6\) 和 \(6\) 个单位. 这些流量分别不超过图中对应网线的容量，且下方两条路线在汇合后通过容量为 \(12\) 的网线，故可行；到达 \(A\) 的总量为 \(3+4+6+6=19\). 上界可以达到，所以最大信息量为 \(19\)，对应选项 D.
\end{solution}

\section{填空题}
\begin{problem}
\(\left(\frac{1}{2} x + 1\right)^{10}\) 的二项展开式中 \(x^{3}\) 的系数为\fillinblank{}.
\end{problem}

\begin{answer}
15
\end{answer}

\begin{solution}
由二项式定理，展开式的通项为 \(\mathrm C_{10}^{k}\left(\frac12x\right)^k\). 要得到 \(x^3\)，必须取 \(k=3\)，所以其系数为
\(\mathrm C_{10}^{3}\left(\frac12\right)^3=120\cdot\frac18=15\).
\end{solution}

\begin{problem}
双曲线 \( \frac{x^{2}}{9}-\frac{y^{2}}{16}=1 \) 的两个焦点为 \(F_{1}\)，\(F_{2}\)，点 \(P\) 在双曲线上，若 \( PF_{1} \perp PF_{2} \)，则点 \(P\) 到 \(x\) 轴的距离为\fillinblank{}.
\end{problem}

\begin{answer}
\(\frac{16}{5}\)
\end{answer}

\begin{solution}
双曲线的焦半距为 \(c=\sqrt{9+16}=5\)，故焦点为 \(F_1(-5,0)\)，\(F_2(5,0)\). 设 \(P=(x,y)\). 由垂直条件
\(\overrightarrow{PF_1}\cdot\overrightarrow{PF_2}=(-5-x,-y)\cdot(5-x,-y)=0\)，得 \(x^2+y^2=25\). 又因为 \(P\) 在双曲线上，满足 \(\frac{x^2}{9}-\frac{y^2}{16}=1\). 将 \(y^2=25-x^2\) 代入，得
\(16x^2-9(25-x^2)=144\)，即 \(25x^2=369\)，所以 \(x^2=\frac{369}{25}\). 于是
\(y^2=25-\frac{369}{25}=\frac{256}{25}\)，从而点 \(P\) 到 \(x\) 轴的距离为 \(|y|=\frac{16}{5}\).
\end{solution}

\begin{problem}
设 \(\{a_{n}\}\) 是公比为 \(q\) 的等比数列，\(S_{n}\) 是它的前 \(n\) 项和. 若 \(\{S_{n}\}\) 是等差数列，则 \(q = \)\fillinblank{}.
\end{problem}

\begin{answer}
1
\end{answer}

\begin{solution}
因为 \(\{S_n\}\) 是等差数列，所以相邻两项的差相等，即 \(S_2-S_1=S_3-S_2\). 由前 \(n\) 项和的定义，\(S_2-S_1=a_2=a_1q\)，\(S_3-S_2=a_3=a_1q^2\)，故
\(a_1q=a_1q^2\). 按等比数列公比的定义，\(a_1\ne0\) 且 \(q\ne0\)，于是可约去 \(a_1q\)，得到 \(q=1\).
\end{solution}

\begin{problem}
圆周上有 \(2n\) 个等分点（\(n > 1\)），以其中三个点为顶点的直角三角形的个数为\fillinblank{}.
\end{problem}

\begin{answer}
\(2n(n-1)\)
\end{answer}

\begin{solution}
若圆周上的三个等分点构成直角三角形，则由圆周角定理，直角所对的弦必须是直径. 反之，任取一个直径，再从其余的 \(2n-2\) 个等分点中任选一个作为第三个顶点，根据直径所对的圆周角为直角，必得到一个直角三角形. \(2n\) 个等分点两两相对，所以直径共有 \(n\) 条；每个直角三角形的直径唯一，不会重复计数. 因此个数为 \(n(2n-2)=2n(n-1)\).
\end{solution}

\section{解答题}
\begin{problem}
已知等差数列前三项为 \(a\)，\(4\)，\(3a\)，前 \(n\) 项的和为 \(S_{n}\)，\(S_{k} = 2550\).

\begin{enumerate}
\item 求 \(a\) 及 \(k\) 的值；
\item 求 \(\displaystyle\lim_{n\to\infty}\left(\frac{1}{S_{1}}+\frac{1}{S_{2}}+\cdots+\frac{1}{S_{n}}\right)\).
\end{enumerate}
\end{problem}

\begin{solution}
\begin{enumerate}
\item 等差数列的第二项是第一项与第三项的平均数，因此 \(4=\frac{a+3a}{2}\)，即 \(a+3a=8\)，解得 \(a=2\). 此时首项为 \(a_1=2\)，公差为 \(d=4-2=2\). 由等差数列前 \(k\) 项和公式，
\(S_k=ka_1+\frac{k(k-1)}{2}d=2k+k(k-1)=k(k+1)\). 代入 \(S_k=2550\)，得 \(k^2+k-2550=0\)，即 \((k-50)(k+51)=0\). 因为 \(k\) 是正整数，所以 \(k=50\)，故 \(a=2\)，\(k=50\).
\item 对任意正整数 \(n\)，同理有 \(S_n=2n+\frac{n(n-1)}{2}\times2=n(n+1)\). 因而
\[
\frac{1}{S_1}+\frac{1}{S_2}+\cdots+\frac{1}{S_n}=\sum_{k=1}^{n}\frac{1}{k(k+1)}=\sum_{k=1}^{n}\left(\frac{1}{k}-\frac{1}{k+1}\right)=1-\frac{1}{n+1}
\]
所以所求极限为 \(\displaystyle\lim_{n\to\infty}\left(1-\frac{1}{n+1}\right)=1\).
\end{enumerate}
\end{solution}

\begin{problem}
如图，在底面是直角梯形的四棱锥 \(S - ABCD\) 中，\(\angle ABC = 90^{\circ}\)，\(SA \perp\) 面 \(ABCD\)，\(SA = AB = BC = 1\)，\(AD = \frac{1}{2}\).

\begin{enumerate}
\item 求四棱锥 \(S - ABCD\) 的体积；
\item 求面 \(SCD\) 与面 \(SBA\) 所成的二面角的正切值.
\end{enumerate}

\bitmapfigure[max width=0.34\linewidth]{img/2001/syllabus_liberal/q18_fig1.png}
\end{problem}

\begin{solution}
\begin{enumerate}
\item 因为底面是直角梯形，且 \(AD\parallel BC\)，所以 \(AB\) 是梯形的高. 底面面积为
\(S_{ABCD}=\frac{AD+BC}{2}\cdot AB=\frac{\frac12+1}{2}\cdot1=\frac34\). 又因为 \(SA\perp\) 面 \(ABCD\)，四棱锥的高为 \(SA=1\)，所以
\(V=\frac13S_{ABCD}\cdot SA=\frac13\cdot\frac34\cdot1=\frac14\).
\item 建立空间直角坐标系，取 \(A=(0,0,0)\)，\(B=(1,0,0)\)，\(C=(1,1,0)\)，\(D=\left(0,\frac12,0\right)\)，\(S=(0,0,1)\). 则平面 \(SBA\) 的一个法向量为 \(\bs{n}_1=(0,1,0)\). 又
\(\overrightarrow{SC}=(1,1,-1)\)，\(\overrightarrow{SD}=\left(0,\frac12,-1\right)\)，所以平面 \(SCD\) 的一个法向量可取 \(\bs{n}_2=\overrightarrow{SC}\times\overrightarrow{SD}=\left(-\frac12,1,\frac12\right)\)，即也可取 \(\bs{n}_2=(-1,2,1)\). 设两平面所成的锐二面角为 \(\theta\)，则
\[
\cos\theta=\frac{|\bs{n}_1\cdot\bs{n}_2|}{|\bs{n}_1||\bs{n}_2|}=\frac{2}{\sqrt{6}}
\]
平面角取锐角时，正切值为正，因此 \(\tan\theta=\frac{\sqrt{1-\cos^2\theta}}{\cos\theta}=\frac{1}{\sqrt{2}}\). 所求二面角的正切值为 \(\frac{1}{\sqrt{2}}\).
\end{enumerate}
\end{solution}

\begin{problem}
已知圆内接四边形 \(ABCD\) 的边长分别为 \(AB = 2\)，\(BC = 6\)，\(CD = DA = 4\)，求四边形 \(ABCD\) 的面积.
\end{problem}

\begin{solution}
连接 \(BD\)，设 \(\angle BAD=A\)，\(\angle BCD=C\). 因为四边形 \(ABCD\) 为圆内接四边形，所以对角互补，\(A+C=180^{\circ}\)，从而 \(\sin A=\sin C\). 四边形面积为两个三角形面积之和，因此
\[
S=\frac12AB\cdot AD\sin A+\frac12BC\cdot CD\sin C=16\sin A
\]
在三角形 \(ABD\) 和 \(CBD\) 中分别使用余弦定理，得
\(BD^2=AB^2+AD^2-2AB\cdot AD\cos A=20-16\cos A\)，
\(BD^2=BC^2+CD^2-2BC\cdot CD\cos C=52-48\cos C\). 又 \(\cos C=-\cos A\)，所以
\(20-16\cos A=52+48\cos A\)，从而 \(\cos A=-\frac12\). 因为 \(0<A<180^{\circ}\)，故 \(\sin A=\frac{\sqrt{3}}{2}\). 因此 \(S=16\cdot\frac{\sqrt{3}}{2}=8\sqrt{3}\).
\end{solution}

\begin{problem}
设抛物线 \(y^{2} = 2px\)（\(p > 0\)）的焦点为 \(F\)，经过点 \(F\) 的直线交抛物线于 \(A\)，\(B\) 两点. 点 \(C\) 在抛物线的准线上，且 \(BC \parallel x\) 轴. 证明直线 \(AC\) 经过原点 \(O\).
\end{problem}

\begin{solution}
抛物线 \(y^2=2px\) 的焦点为 \(F\left(\frac p2,0\right)\)，准线为 \(x=-\frac p2\). 设抛物线上点 \(A\)，\(B\) 的参数分别为 \(t_1\)，\(t_2\)，则
\(A=\left(\frac p2t_1^2,pt_1\right)\)，\(B=\left(\frac p2t_2^2,pt_2\right)\). 参数为 \(t_1\)，\(t_2\) 的两点所成弦的方程为
\(2x-(t_1+t_2)y+pt_1t_2=0\). 直线 \(AB\) 经过焦点，代入 \(F\left(\frac p2,0\right)\) 得 \(p+pt_1t_2=0\)，所以 \(t_1t_2=-1\).

由 \(BC\parallel x\) 轴，点 \(C\) 与点 \(B\) 的纵坐标相同；又 \(C\) 在准线上，所以 \(C=\left(-\frac p2,pt_2\right)\). 因为 \(t_1t_2=-1\)，所以 \(t_1\ne0\)，从而
\(k_{OA}=\frac{pt_1}{\frac p2t_1^2}=\frac{2}{t_1}\)，而
\(k_{OC}=\frac{pt_2}{-\frac p2}=-2t_2=\frac{2}{t_1}\). 两条直线均经过 \(O\) 且斜率相同，因此 \(O\)，\(A\)，\(C\) 三点共线，直线 \(AC\) 经过原点 \(O\).
\end{solution}

\begin{problem}
设计一幅宣传画，要求画面面积为 \(4840 \mathrm{~cm}^{2}\)，画面的宽与高的比为 \(\lambda\)（\(\lambda < 1\)），画面的上、下各留 \(8 \mathrm{~cm}\) 空白，左、右各留 \(5 \mathrm{~cm}\) 空白. 怎样确定画面的高与宽尺寸，能使宣传画所用纸张面积最小？
\end{problem}

\begin{solution}
设画面高为 \(h\mathrm{~cm}\)，宽为 \(w\mathrm{~cm}\)，则 \(wh=4840\)，且 \(\lambda=\frac{w}{h}<1\). 纸张的高为 \(h+16\)，宽为 \(w+10\)，所以纸张面积为
\[
S=(h+16)(w+10)=4840+16w+10h+160=5000+16w+10h
\]
因为 \(w\)，\(h>0\)，由均值不等式
\(16w+10h\geqslant2\sqrt{160wh}=2\sqrt{160\times4840}=1760\)，等号成立的条件是 \(16w=10h\)，即 \(\frac{w}{h}=\frac58\). 此时宽高比满足 \(\lambda=\frac58<1\)，符合题设. 由 \(w=\frac58h\) 和 \(wh=4840\) 得 \(\frac58h^2=4840\)，所以 \(h=88\)，\(w=55\). 此时纸张面积达到最小值
\(S_{\min}=5000+1760=6760\mathrm{~cm}^{2}\). 因而画面高为 \(88\mathrm{~cm}\)，宽为 \(55\mathrm{~cm}\) 时最优.
\end{solution}

\begin{problem}
设 \(f(x)\) 是定义在 \(\R\) 上的偶函数，其图象关于直线 \(x = 1\) 对称，对任意 \(x_{1}\)，\(x_{2} \in \left[0, \frac{1}{2}\right]\)，都有 \(f(x_{1} + x_{2}) = f(x_{1}) \cdot f(x_{2})\).

\begin{enumerate}
\item 设 \(f(1)=2\)，求 \(f\left(\frac{1}{2}\right)\)，\(f\left(\frac{1}{4}\right)\)；
\item 证明：\(f(x)\) 是周期函数.
\end{enumerate}
\end{problem}

\begin{solution}
\begin{enumerate}
\item 对任意 \(x\in[0,1]\)，令 \(x_1=x_2=\frac{x}{2}\)，则 \(x_1\)，\(x_2\in\left[0,\frac12\right]\)，由题设得
\(f(x)=f\left(\frac{x}{2}\right)^2\geqslant0\). 特别地，令 \(x_1=x_2=\frac12\)，得到
\(f(1)=f\left(\frac12\right)^2\). 已知 \(f(1)=2\)，且 \(f\left(\frac12\right)\geqslant0\)，故 \(f\left(\frac12\right)=\sqrt{2}=2^{\frac12}\). 再令 \(x_1=x_2=\frac14\)，得到
\(f\left(\frac12\right)=f\left(\frac14\right)^2\). 同样由非负性取正根，得 \(f\left(\frac14\right)=\sqrt{\sqrt{2}}=2^{\frac14}\).
\item 函数图象关于直线 \(x=1\) 对称，等价于对任意 \(x\in\R\) 有 \(f(x)=f(2-x)\). 又因为 \(f\) 是偶函数，所以 \(f(-x)=f(x)\). 在对称关系中以 \(-x\) 代替 \(x\)，得 \(f(-x)=f(2+x)\). 合并两式可得 \(f(x)=f(x+2)\)，对任意 \(x\in\R\) 成立. 因而 \(f(x)\) 是周期函数，且 \(2\) 是它的一个周期.
\end{enumerate}
\end{solution}
